\section{Gibbs measures with symmetries}

Throughout this section $S$ is the site set, $(E,\mathcal{E})$ the state space,
$\Cfg = E^S$ the configuration space, and $T$ the group of transformations
$\tau = (\tau_*; \tau_i, i \in S)$ of $\Cfg$ of Georgii \S5.1: $\tau_*$ is a permutation of $S$,
each $\tau_i$ a bimeasurable bijection of $E$, and $(\tau\omega)_{\tau_* i} = \tau_i(\omega_i)$.
A transformation $\tau$ is a \textbf{symmetry} of a specification $\gamma$ when
$\tau(\gamma) = \gamma$, where $\tau(\gamma)$ is the image specification of Georgii (5.4).

\subsection*{Transport of Gibbs measures}

\begin{remark}[Georgii (5.10)]
    \label{rmk:sym-transport-gibbs}
    \uses{def:spec, def:gibbs-meas}
    \lean{Specification.map_mem_GP,
          Specification.IsInvariant.map_mem_GP,
          Specification.IsInvariant.map_isGibbsMeasure}
    \leanok

    Let $\gamma$ be a specification and $\tau \in T$. If $\mu \in \Gibbs{\gamma}$ then
    $\tau(\mu) \in \mathcal{G}(\tau(\gamma))$. In particular $\Gibbs{\gamma}$ is invariant under
    every symmetry of $\gamma$.
\end{remark}
\begin{proof}
    \leanok

    For $\Lambda \in \Fin$, using (5.4) and then the DLR equation for $\mu$ at the volume
    $\tau_*^{-1}\Lambda$,
    \[
      \tau(\mu)\,\tau(\gamma)_\Lambda
        = \int \mu(\mathrm{d}\omega)\, \tau(\gamma)_\Lambda(\,\cdot\, | \tau\omega)
        = \int \mu(\mathrm{d}\omega)\, \gamma_{\tau_*^{-1}\Lambda}(\tau^{-1}\,\cdot\,|\omega)
        = \tau\bigl(\mu\gamma_{\tau_*^{-1}\Lambda}\bigr) = \tau(\mu).
    \]
\end{proof}

\begin{corollary}[Georgii (5.11)]
    \label{cor:sym-unique-invariant}
    \uses{rmk:sym-transport-gibbs}
    \lean{Specification.IsInvariant.measurePreserving_of_GP_eq_singleton}
    \leanok

    If $\Gibbs{\gamma} = \set{\mu}$ then $\mu$ is preserved by every symmetry of $\gamma$.
\end{corollary}
\begin{proof}
    \leanok

    By Remark \ref{rmk:sym-transport-gibbs}, $\tau(\mu) \in \Gibbs{\gamma} = \set{\mu}$, so
    $\tau(\mu) = \mu$.
\end{proof}

\subsection*{Invariant random fields}

\begin{definition}[Georgii (5.12), (5.13)]
    \label{def:sym-invariant-fields}
    \uses{def:gibbs-meas}

    For $I \subseteq T$ let
    \[
      \Prob_I(\Cfg,\mathcal{F}) = \set{\mu \in \Prob(\Cfg,\mathcal{F}) :
        \tau(\mu) = \mu \text{ for all } \tau \in I}
    \]
    be the set of $I$-invariant random fields, and
    $\mathcal{G}_I(\gamma) = \Gibbs{\gamma} \cap \Prob_I(\Cfg,\mathcal{F})$ the set of
    $I$-invariant Gibbs measures. Since $\Prob_I = \Prob_{[I]}$ for the group $[I]$ generated by
    $I$, nothing is lost by taking $I$ to be a group.
\end{definition}

\begin{lemma}[Georgii (5.12), (5.13): $\Prob_I$ and $\mathcal{G}_I(\gamma)$ are closed]
    \label{lem:sym-closed-invariant}
    \uses{def:sym-invariant-fields, def:cylinder-event}
    \lean{MeasureTheory.GibbsMeasure.isClosed_setOf_measurePreserving,
          MeasureTheory.GibbsMeasure.isClosed_setOf_forall_measurePreserving,
          MeasureTheory.GibbsMeasure.isClosed_setOf_mem_GP_and_measurePreserving}
    \leanok

    $\Prob_I(\Cfg,\mathcal{F})$ is closed in the topology $\mathcal{L}$ of local convergence, and
    for a quasilocal specification $\gamma$ so is $\mathcal{G}_I(\gamma)$.
\end{lemma}
\begin{proof}
    \leanok

    The preimage of a local event under $\tau \in T$ is again a local event, so
    $\mu \mapsto \tau(\mu)$ is $\mathcal{L}$-continuous and $\Prob_{\set{\tau}}$, the equaliser of
    two continuous maps, is closed; $\Prob_I$ is the intersection over $\tau \in I$. Intersecting
    with the $\mathcal{L}$-closed set $\Gibbs{\gamma}$ of Georgii (4.17) gives the second claim.
\end{proof}

\subsection*{Averaging machinery}

Both existence proofs average a family of probability measures, over transformations in one case
and over volumes in the other.

\begin{definition}[Uniform average of a finite family of measures]
    \label{def:sym-uniform-average}
    \lean{MeasureTheory.uniformAverage,
          MeasureTheory.isProbabilityMeasure_uniformAverage,
          MeasureTheory.uniformAverage_real_apply}
    \leanok

    For a family $(m_i)_{i \in \iota}$ of measures and a finite $F \subseteq \iota$ set
    $\mathrm{avg}_F m = |F|^{-1}\sum_{i \in F} m_i$. If every $m_i$ is a probability measure and
    $F \neq \emptyset$, so is $\mathrm{avg}_F m$.
\end{definition}

\begin{lemma}[Comparison of uniform averages]
    \label{lem:sym-uniform-average-estimate}
    \uses{def:sym-uniform-average}
    \lean{MeasureTheory.abs_uniformAverage_real_sub_le,
          MeasureTheory.abs_uniformAverage_real_sub_le_of_card_eq}
    \leanok

    For non-empty finite $F, F'$, probability measures $m_i$, and every event $A$,
    \[
      \bigl|\mathrm{avg}_F m(A) - \mathrm{avg}_{F'} m(A)\bigr|
        \le \frac{|F \mathbin{\triangle} F'|}{|F|} + \left|\frac{|F'|}{|F|} - 1\right| ,
    \]
    and in particular $\le |F \mathbin{\triangle} F'| / |F|$ when $|F| = |F'|$.
\end{lemma}
\begin{proof}
    \leanok

    Write $g(i) = m_i(A) \in [0,1]$. Then
    $|F|^{-1}\sum_{F} g - |F'|^{-1}\sum_{F'} g
      = |F|^{-1}\bigl(\sum_{F \setminus F'} g - \sum_{F' \setminus F} g\bigr)
        + (|F|^{-1} - |F'|^{-1})\sum_{F'} g$.
    Bound the first bracket by $|F \mathbin{\triangle} F'|$ using $0 \le g \le 1$, and the second
    summand by $\bigl||F'|/|F| - 1\bigr|$.
\end{proof}

\begin{lemma}[Finite convex combinations of Gibbs measures]
    \label{lem:sym-convex-combination}
    \uses{def:gibbs-meas, def:sym-uniform-average}
    \lean{MeasureTheory.mem_GP_finset_sum_smul, MeasureTheory.mem_GP_uniformAverage}
    \leanok

    A finite linear combination $\sum_{i \in F} w_i\, m_i$ of Gibbs measures for $\gamma$ which
    is again a probability measure is a Gibbs measure for $\gamma$; in particular
    $\mathrm{avg}_F m \in \Gibbs{\gamma}$.
\end{lemma}
\begin{proof}
    \leanok

    $\mathrm{bind}$ against the kernel $\gamma_\Lambda$ commutes with finite sums and with
    scalar multiplication of measures, so $\mu\gamma_\Lambda = \sum_i w_i\, m_i \gamma_\Lambda
    = \sum_i w_i\, m_i = \mu$ for every $\Lambda \in \Fin$.
\end{proof}

\begin{definition}[Average over a family of transformations, Georgii (5.15)]
    \label{def:sym-trans-average}
    \uses{def:sym-uniform-average}
    \lean{MeasureTheory.GibbsMeasure.transAverage,
          MeasureTheory.GibbsMeasure.isProbabilityMeasure_transAverage}
    \leanok

    For an action $\Phi : A \to T$ of an index type $A$ by transformations, a random field $\nu$
    and a non-empty finite $F \subseteq A$, put
    $\mathrm{avg}_F \Phi(\nu) = |F|^{-1}\sum_{a \in F} \Phi(a)(\nu)$; it is a probability measure.
\end{definition}

\begin{lemma}[Translating the index set, Georgii (5.15)(ii)]
    \label{lem:sym-map-trans-average}
    \uses{def:sym-trans-average}
    \lean{MeasureTheory.GibbsMeasure.map_transAverage}
    \leanok

    If $A$ is an abelian group and $\Phi(x + y) = \Phi(x) \circ \Phi(y)$, then
    $\Phi(b)\bigl(\mathrm{avg}_F \Phi(\nu)\bigr) = \mathrm{avg}_{b + F} \Phi(\nu)$ for all
    $b \in A$.
\end{lemma}
\begin{proof}
    \leanok

    Push-forward commutes with finite sums and scalars, and
    $\Phi(b)\bigl(\Phi(a)(\nu)\bigr) = \Phi(b+a)(\nu)$ by the homomorphism property; reindex the
    sum by $a \mapsto b + a$, a bijection of $F$ onto $b + F$.
\end{proof}

\begin{theorem}[Every abelian group is amenable: F\o lner sets]
    \label{thm:sym-foelner}
    \lean{Finset.transDist, AddCommGroup.exists_finset_transDist_le,
          AddCommGroup.exists_foelner_seq, AddCommGroup.exists_isAddFoelner_count,
          AddCommGroup.exists_invariant_finitelyAdditive}
    \leanok

    Let $G$ be an abelian group, $s \subseteq G$ finite and $\varepsilon > 0$. Then there is a
    non-empty finite $F \subseteq G$, contained in the subgroup generated by $s$, with
    $|(g + F) \mathbin{\triangle} F| \le \varepsilon |F|$ for all $g \in s$. For countable $G$
    these assemble into a F\o lner sequence, and hence into a translation-invariant finitely
    additive probability measure on $G$.
\end{theorem}
\begin{proof}
    \leanok

    Induction on $s$. For $s = \emptyset$ take $F = \set{0}$. To add a generator $a$ to a set $F$
    contained in the subgroup $H$ generated by the previous generators, replace $F$ by
    $\bigcup_{k < M} (k a + F)$, where $M$ is the order of $a$ in $G/H$ if that is finite and is
    chosen large otherwise. The slices lie in distinct cosets, so cardinalities add: the ratios
    already achieved are unchanged and the new one is $O(1/M)$.
\end{proof}

\subsection*{First approach: symmetrisation, Georgii (5.15)}

\begin{theorem}[Georgii (5.15)]
    \label{thm:sym-5-15}
    \uses{def:sym-invariant-fields}

    Let $I_0, I_1$ be subgroups of $T$ with $\tau_1 \circ I_0 = I_0 \circ \tau_1$ for all
    $\tau_1 \in I_1$, and let $\gamma$ be an $I_1 \circ I_0$-invariant specification. Suppose
    that either
    \begin{enumerate}
        \item there is a topology making $I_1$ a compact topological group for which the
              evaluation map $\mathbf{e} : (\tau,\omega) \mapsto \tau\omega$ is measurable; or
        \item $\mathcal{G}_{I_0}(\gamma)$ is compact in the topology of local convergence, and
              for all $\tau_1, \tau_2 \in I_1$ there is $\tau_0 \in I_0$ with
              $\tau_1 \circ \tau_2 = \tau_2 \circ \tau_1 \circ \tau_0$.
    \end{enumerate}
    Then $\mathcal{G}_{I_0}(\gamma) \neq \emptyset$ implies
    $\mathcal{G}_{I_1 \circ I_0}(\gamma) \neq \emptyset$.
\end{theorem}

\begin{theorem}[Georgii (5.15)(ii) with $I_0 = \set{\mathrm{id}}$]
    \label{thm:sym-5-15-abelian}
    \uses{thm:sym-5-15, lem:sym-convex-combination, lem:sym-map-trans-average, thm:sym-foelner, rmk:sym-transport-gibbs, lem:sym-uniform-average-estimate}
    \lean{MeasureTheory.GibbsMeasure.exists_mem_GP_and_forall_measurePreserving_of_isCompact}
    \leanok

    Let $A$ be an abelian group acting on $\Cfg$ by a homomorphism $\Phi : A \to T$, not assumed
    injective, so $\Phi(x+y) = \Phi(x) \circ \Phi(y)$, and let every $\Phi(x)$ be a symmetry of
    $\gamma$. If $\Gibbs{\gamma}$ is non-empty and compact in the topology of local convergence,
    then it contains a measure invariant under every $\Phi(x)$.
\end{theorem}
\begin{proof}
    \leanok

    Georgii's Markov--Kakutani argument. Pick $\nu \in \Gibbs{\gamma}$. By
    Theorem \ref{thm:sym-foelner} choose, for each finite set of directions and each accuracy,
    a F\o lner set $F$; the averages $\mu_F = \mathrm{avg}_F \Phi(\nu)$ are Gibbs measures by
    Remark \ref{rmk:sym-transport-gibbs} and Lemma \ref{lem:sym-convex-combination}. Compactness
    of $\Gibbs{\gamma}$ gives a cluster point $\mu \in \Gibbs{\gamma}$ along the net of these
    averages. For a local event $B$ and $x \in A$, Lemma \ref{lem:sym-map-trans-average} turns
    $\mu_F(\Phi(x)^{-1}B)$ into the average over $x + F$, and
    Lemma \ref{lem:sym-uniform-average-estimate} bounds the difference from $\mu_F(B)$ by
    $|(x+F) \mathbin{\triangle} F| / |F| \to 0$. Hence $\mu(\Phi(x)^{-1}B) = \mu(B)$ on local
    events, which separate probability measures.
\end{proof}

\begin{corollary}[Georgii (5.16), abelian branch]
    \label{cor:sym-5-16-abelian}
    \uses{thm:sym-5-15-abelian}
    \lean{MeasureTheory.GibbsMeasure.exists_mem_GP_and_forall_measurePreserving_of_commute}
    \leanok

    Let $I$ be an abelian subgroup of $T$ and $\gamma$ an $I$-invariant specification with
    $\Gibbs{\gamma} \neq \emptyset$. If $\Gibbs{\gamma}$ is compact then
    $\mathcal{G}_I(\gamma) \neq \emptyset$.
\end{corollary}
\begin{proof}
    \leanok

    Apply Theorem \ref{thm:sym-5-15-abelian} to $A = I$ viewed additively and $\Phi$ the
    inclusion, using that the elements of $I$ commute pairwise.
\end{proof}

\begin{definition}[Average over a group action against a weight]
    \label{def:sym-group-average}
    \lean{MeasureTheory.GibbsMeasure.groupAverage,
          MeasureTheory.GibbsMeasure.measurable_map_toFun,
          MeasureTheory.GibbsMeasure.isProbabilityMeasure_groupAverage}
    \leanok

    For a group $H$ acting by $\Phi : H \to T$, a random field $\nu$ and a measure $w$ on $H$,
    put $\mathrm{avg}_w \Phi(\nu) = \int_H \Phi(g)(\nu)\, w(\mathrm{d}g)$. If the evaluation map
    $\mathbf{e} : (g,\omega) \mapsto \Phi(g)\omega$ is measurable then $g \mapsto \Phi(g)(\nu)$ is
    a measurable map into the space of measures, so the integral is defined; if $\nu$ and $w$ are
    probability measures then so is $\mathrm{avg}_w \Phi(\nu)$.
\end{definition}

\begin{lemma}[The weighted average is Gibbs and invariant]
    \label{lem:sym-group-average-invariant}
    \uses{def:sym-group-average, rmk:sym-transport-gibbs, def:gibbs-meas}
    \lean{MeasureTheory.GibbsMeasure.isGibbsMeasure_groupAverage,
          MeasureTheory.GibbsMeasure.map_groupAverage_eq}
    \leanok

    If every $\Phi(g)$ is a symmetry of $\gamma$ and $\nu \in \Gibbs{\gamma}$, then
    $\mathrm{avg}_w \Phi(\nu) \in \Gibbs{\gamma}$. If moreover $\Phi$ is a homomorphism and $w$
    is left invariant, then $\Phi(h)\bigl(\mathrm{avg}_w \Phi(\nu)\bigr) = \mathrm{avg}_w
    \Phi(\nu)$ for every $h \in H$.
\end{lemma}
\begin{proof}
    \leanok

    Each $\Phi(g)(\nu)$ is Gibbs by Remark \ref{rmk:sym-transport-gibbs}, and
    $\mathrm{bind}$ is associative, so $\bigl(\mathrm{avg}_w \Phi(\nu)\bigr)\gamma_\Lambda
    = \int \Phi(g)(\nu)\gamma_\Lambda\, w(\mathrm{d}g) = \mathrm{avg}_w \Phi(\nu)$. For the
    second claim, $\Phi(h)\bigl(\Phi(g)(\nu)\bigr) = \Phi(hg)(\nu)$, so pushing the average
    forward by $\Phi(h)$ replaces $w$ by its image under $g \mapsto hg$, which is $w$ again.
\end{proof}

\begin{theorem}[Georgii (5.15)(i), at the hypothesis its proof uses]
    \label{thm:sym-5-15-weight}
    \uses{thm:sym-5-15, lem:sym-group-average-invariant}
    \lean{MeasureTheory.GibbsMeasure.exists_isGibbsMeasure_and_forall_map_eq_of_invariantWeight,
          MeasureTheory.GibbsMeasure.exists_mem_GP_and_forall_measurePreserving_of_invariantWeight}
    \leanok

    Let a group $H$ with measurable multiplication act on $\Cfg$ by symmetries $\Phi(g)$ of
    $\gamma$ with measurable evaluation map $(g,\omega) \mapsto \Phi(g)\omega$, and suppose $H$
    carries a left-invariant probability measure $w$ --- in place of a compact topology, which
    enters Georgii's hypothesis (i) only through its normalised Haar measure
    (Corollary \ref{cor:sym-5-16-compact}). If $\Gibbs{\gamma} \neq \emptyset$ then
    $\Gibbs{\gamma}$ contains an $H$-invariant measure, with no compactness assumption on
    $\Gibbs{\gamma}$.
\end{theorem}
\begin{proof}
    \leanok

    Take $\nu \in \Gibbs{\gamma}$ and set $\mu = \int_H \Phi(g)(\nu)\, w(\mathrm{d}g)$, which is
    a probability measure by Definition \ref{def:sym-group-average}. By
    Lemma \ref{lem:sym-group-average-invariant} it lies in $\Gibbs{\gamma}$ and is
    $H$-invariant.
\end{proof}

\begin{corollary}[Georgii (5.16), compact branch]
    \label{cor:sym-5-16-compact}
    \uses{thm:sym-5-15-weight}
    \lean{MeasureTheory.GibbsMeasure.haarProb,
          MeasureTheory.GibbsMeasure.map_mul_haarProb,
          MeasureTheory.GibbsMeasure.exists_mem_GP_and_forall_measurePreserving_of_compactGroup}
    \leanok

    Let $H$ be a compact topological group with Borel structure acting on $\Cfg$ by symmetries of
    $\gamma$ with measurable evaluation map. If $\Gibbs{\gamma} \neq \emptyset$ then
    $\mathcal{G}_H(\gamma) \neq \emptyset$.
\end{corollary}
\begin{proof}
    \leanok

    The normalised Haar measure of a compact group is a left-invariant probability measure; apply
    Theorem \ref{thm:sym-5-15-weight} with $w$ equal to it.
\end{proof}

\subsection*{Second approach: construction from the specification, Georgii (5.18)--(5.19)}

\begin{definition}[Averaged Gibbs distributions, Georgii (5.18)]
    \label{def:sym-average}
    \uses{def:spec, def:markov-ker}
    \lean{Specification.average, Specification.isProbabilityMeasure_average,
          Specification.average_apply}
    \leanok

    For a specification $\gamma$, a random field $\nu$ and a non-empty finite family
    $\mathcal{R} \subseteq \Fin$ of volumes,
    \[
      \mathrm{avg}_{\mathcal{R}}(\nu,\gamma)
        = |\mathcal{R}|^{-1} \sum_{\Lambda \in \mathcal{R}} \nu\gamma_\Lambda ;
    \]
    it is a probability measure.
\end{definition}

\begin{lemma}[Transport and comparison of averages, Georgii (5.5), (5.18)]
    \label{lem:sym-average-transport}
    \uses{def:sym-average, lem:sym-uniform-average-estimate}
    \lean{Specification.map_average, Specification.abs_average_real_sub_le,
          Specification.abs_average_real_sub_le_of_card_eq,
          Specification.abs_average_real_sub_le_of_subset}
    \leanok

    Let $\tau$ be a symmetry of $\gamma$ preserving $\nu$. Then
    $\tau\bigl(\mathrm{avg}_{\mathcal{R}}(\nu,\gamma)\bigr)
      = \mathrm{avg}_{\tau_*\mathcal{R}}(\nu,\gamma)$, where
    $\tau_*\mathcal{R} = \set{\tau_*\Lambda : \Lambda \in \mathcal{R}}$. Moreover, for non-empty
    $\mathcal{R}, \mathcal{R}'$ of equal cardinality and every event $A$,
    $\bigl|\mathrm{avg}_{\mathcal{R}}(A) - \mathrm{avg}_{\mathcal{R}'}(A)\bigr|
      \le |\mathcal{R} \mathbin{\triangle} \mathcal{R}'| / |\mathcal{R}|$; and for
    $\mathcal{R}' \subseteq \mathcal{R}$ the bound
    $2\bigl(1 - |\mathcal{R}'|/|\mathcal{R}|\bigr)$ holds without the cardinality assumption.
\end{lemma}
\begin{proof}
    \leanok

    For a single volume, $\tau(\nu\gamma_\Lambda) = \nu\gamma_{\tau_*\Lambda}$: push
    $\mathrm{bind}$ through $\tau$, use $\tau(\gamma_\Lambda(\cdot|\omega)) =
    \gamma_{\tau_*\Lambda}(\cdot|\tau\omega)$ from (5.5) and then $\tau(\nu) = \nu$. Summing and
    renormalising gives the first claim, since $\tau_*$ is injective on volumes. The estimates
    are Lemma \ref{lem:sym-uniform-average-estimate} with index set $\Fin$; the last one uses
    $|\mathcal{R} \mathbin{\triangle} \mathcal{R}'| / |\mathcal{R}| =
    1 - |\mathcal{R}'|/|\mathcal{R}|$ for $\mathcal{R}' \subseteq \mathcal{R}$.
\end{proof}

\begin{lemma}[Averages over large volumes are fixed by $\gamma_\Lambda$, Georgii (4.18)]
    \label{lem:sym-consistency-average}
    \uses{def:sym-average, def:spec}
    \lean{Specification.bind_average_of_subset}
    \leanok

    If $\Lambda \subseteq \Lambda'$ for every $\Lambda' \in \mathcal{R}$ then
    $\mathrm{avg}_{\mathcal{R}}(\nu,\gamma)\, \gamma_\Lambda
      = \mathrm{avg}_{\mathcal{R}}(\nu,\gamma)$.
\end{lemma}
\begin{proof}
    \leanok

    Termwise: $\nu\gamma_{\Lambda'}\gamma_\Lambda = \nu\gamma_{\Lambda'}$ by the consistency of
    $\gamma$ for $\Lambda \subseteq \Lambda'$.
\end{proof}

\begin{proposition}[Georgii (5.18)]
    \label{prop:sym-5-18}
    \uses{def:sym-average, lem:sym-average-transport, def:sym-invariant-fields}
    \lean{MeasureTheory.GibbsMeasure.measurePreserving_of_mapClusterPt_average,
          MeasureTheory.GibbsMeasure.measurePreserving_of_mapClusterPt_average_net}
    \leanok

    Let $\tau \in T$ be fixed and let $D$ be a directed set. Suppose given, for each
    $\alpha \in D$, a $\tau$-invariant specification $\gamma^\alpha$, a $\tau$-invariant random
    field $\nu_\alpha$ and a non-empty finite $\mathcal{R}_\alpha \subseteq \Fin$ with
    \[
      |\set{\tau_*\Lambda : \Lambda \in \mathcal{R}_\alpha} \mathbin{\triangle}
        \mathcal{R}_\alpha| \, / \, |\mathcal{R}_\alpha| \longrightarrow 0 .
    \]
    Then every cluster point of
    $\mu_\alpha = |\mathcal{R}_\alpha|^{-1}\sum_{\Lambda \in \mathcal{R}_\alpha}
    \nu_\alpha \gamma^\alpha_\Lambda$ in the topology of local convergence is $\tau$-invariant.
\end{proposition}
\begin{proof}
    \leanok

    Let $\mu$ be a cluster point, realised along an ultrafilter $U$. For a local event $A$, both
    $\mu_\alpha(A)$ and $\mu_\alpha(\tau^{-1}A)$ converge along $U$, the latter because
    $\tau^{-1}A$ is again a local event. By Lemma \ref{lem:sym-average-transport},
    $\mu_\alpha(\tau^{-1}A)$ is the average over $\tau_*\mathcal{R}_\alpha$ evaluated at $A$, and
    $\tau_*\mathcal{R}_\alpha$ has the same cardinality as $\mathcal{R}_\alpha$; hence the same
    lemma bounds $|\mu_\alpha(\tau^{-1}A) - \mu_\alpha(A)|$ by
    $|\tau_*\mathcal{R}_\alpha \mathbin{\triangle} \mathcal{R}_\alpha| / |\mathcal{R}_\alpha|
    \to 0$. So $\mu(\tau^{-1}A) = \mu(A)$ for all local events $A$, which determine $\mu$.
\end{proof}

\begin{lemma}[Local equicontinuity over a finite state space, cf.\ Georgii (4.11)(2)]
    \label{lem:sym-equicontinuous-finite}
    \uses{def:cylinder-event}
    \lean{MeasureTheory.GibbsMeasure.locallyEquicontinuous_of_finite}
    \leanok

    If $E$ is finite and non-empty, every family of random fields on $E^S$ is locally
    equicontinuous.
\end{lemma}
\begin{proof}
    \leanok

    A finite volume $\Lambda$ carries only finitely many events in
    $\mathcal{F}_\Lambda$, since $E^\Lambda$ is finite. So an antitone sequence of such events
    with empty intersection is eventually empty, and the equicontinuity condition holds trivially
    from that point on.
\end{proof}

\begin{theorem}[Georgii (5.19)]
    \label{thm:sym-5-19}
    \uses{prop:sym-5-18, lem:sym-consistency-average, def:sym-invariant-fields, def:gibbs-meas}
    \lean{MeasureTheory.GibbsMeasure.mem_GP_and_forall_measurePreserving_of_mapClusterPt_average,
          MeasureTheory.GibbsMeasure.exists_mem_GP_and_forall_measurePreserving}
    \leanok

    Let $(E,\mathcal{E})$ be standard Borel, $I \subseteq T$, and $\gamma$ a quasilocal
    specification, not itself assumed $I$-invariant. Suppose there exist
    \begin{enumerate}
        \item a net $(\gamma^\alpha)_{\alpha \in D}$ of $I$-invariant specifications with
              $\gamma^\alpha \to \gamma$;
        \item a net $(\mathcal{R}_\alpha)$ of non-empty finite subsets of $\Fin$ with
              $\Lambda_\alpha = \bigcap_{\Lambda \in \mathcal{R}_\alpha} \Lambda \uparrow S$ and
              $|\set{\tau_*\Lambda : \Lambda \in \mathcal{R}_\alpha} \mathbin{\triangle}
              \mathcal{R}_\alpha| / |\mathcal{R}_\alpha| \to 0$ for all $\tau \in I$;
        \item a net $(\nu_\alpha)$ in $\Prob_I(\Cfg,\mathcal{F})$ such that the averages
              $\mu_\alpha = |\mathcal{R}_\alpha|^{-1}\sum_{\Lambda \in \mathcal{R}_\alpha}
              \nu_\alpha\gamma^\alpha_\Lambda$ are locally equicontinuous.
    \end{enumerate}
    Then $\mathcal{G}_I(\gamma)$ contains a cluster point of $(\mu_\alpha)$, and is in particular
    non-empty.
\end{theorem}
\begin{proof}
    \leanok

    Local equicontinuity plus the standard Borel assumption give a cluster point $\mu$ of
    $(\mu_\alpha)$ by Georgii's Proposition (4.9). By
    Lemma \ref{lem:sym-consistency-average}, $\mu_\alpha \gamma^\alpha_{\Lambda_\alpha}
    = \mu_\alpha$, and $\Lambda_\alpha \uparrow S$ with $\gamma^\alpha \to \gamma$, so
    Theorem (4.17) puts $\mu \in \Gibbs{\gamma}$. Proposition \ref{prop:sym-5-18}, applied to
    each $\tau \in I$, makes $\mu$ $I$-invariant.
\end{proof}

\begin{example}[Georgii (5.17)(1), (5.20)(1): configurational boundary conditions on
                $\mathbb{Z}^d$]
    \label{ex:sym-5-20-1}
    \uses{thm:sym-5-19, prop:sym-5-18, lem:sym-equicontinuous-finite, lem:sym-average-transport}
    \lean{MeasureTheory.GibbsMeasure.mem_GP_and_measurePreserving_shift_of_mapClusterPt_average_cubeTranslates,
          MeasureTheory.GibbsMeasure.exists_mem_GP_forall_measurePreserving_shift,
          MeasureTheory.GibbsMeasure.exists_latticeIsing_mem_GP_forall_measurePreserving_shift}
    \leanok

    Let $S = \mathbb{Z}^d$, $\Lambda_N = S \cap [-N,N]^d$, and let $\gamma$ be a quasilocal
    shift-invariant specification. For a shift-invariant boundary condition $\nu$, every cluster
    point of
    $\mu_N = |\Lambda_N|^{-1}\sum_{i \in \Lambda_N} \nu\gamma_{\Lambda_N + i}$
    is a shift-invariant Gibbs measure. If $E$ is finite, non-empty and has measurable
    singletons, such cluster points exist, so
    $\mathcal{G}_\Theta(\gamma) \neq \emptyset$ for the shift group $\Theta$; in particular the
    Ising specification on $\mathbb{Z}^d$ has a shift-invariant Gibbs measure for all coupling,
    field and inverse temperature.
\end{example}
\begin{proof}
    \leanok

    Take $D = \mathbb{N}$, $\gamma^N = \gamma$, $\nu_N = \nu$ and
    $\mathcal{R}_N = \set{\Lambda_N + i : i \in \Lambda_{k(N)}}$, with $k(N) = N$ for the
    first claim. Then
    $\set{\theta_{j*}\Lambda : \Lambda \in \mathcal{R}_N} = \set{\Lambda_N + i :
    i \in \Lambda_{k(N)} + j}$, so the symmetric-difference ratio equals
    $|(\Lambda_{k(N)} + j) \mathbin{\triangle} \Lambda_{k(N)}| / |\Lambda_{k(N)}| \to 0$, and
    Proposition \ref{prop:sym-5-18} gives shift invariance. Passing to $k(N) = N - \sqrt{N}$, one has
    $k(N)/N \to 1$ and $N - k(N) \to \infty$, so $\Lambda_{\sqrt N} \subseteq
    \bigcap \mathcal{R}_N \uparrow S$ while the last estimate of
    Lemma \ref{lem:sym-average-transport} shows the modified averages have the same cluster
    points as the unmodified ones; Theorem \ref{thm:sym-5-19} then applies. Over a finite $E$ the
    equicontinuity hypothesis is automatic by Lemma \ref{lem:sym-equicontinuous-finite}, and
    $\nu = \delta_\omega$ with $\omega$ constant is a shift-invariant boundary condition.
\end{proof}

\begin{remark}[Georgii (5.17)(2)--(4), (5.20)(2)--(3)]
    \label{rmk:sym-5-20-others}
    \uses{thm:sym-5-19, thm:sym-5-15}

    Further groups of symmetries on $\mathbb{Z}^d$ are treated the same way. The reflection and
    lattice-rotation groups $R \circ \Theta$ and the spin-rotation groups
    $U \circ R \circ \Theta$ of (5.17)(2)--(3), and the even/odd sublattice flips of (5.17)(4),
    are instances of Theorem \ref{thm:sym-5-15} with $I_0 = \Theta$ rather than of its
    $I_0 = \set{\mathrm{id}}$ case. On the (5.20) side, Georgii treats free boundary conditions
    with the truncated potentials $\Phi^{\Lambda_N}$ and periodic boundary conditions with the
    periodic modifications $\tilde\Phi^{\Delta_N}$; the latter uses the form of
    Proposition \ref{prop:sym-5-18} with a varying transformation $\tau_\alpha$ subject only to
    $\|f \circ \tau_\alpha - f \circ \tau\| \to 0$ for $f \in \mathcal{L}$.
\end{remark}

\subsection*{Broken symmetries}

\begin{definition}[Georgii (5.21)]
    \label{def:sym-broken}
    \uses{def:gibbs-meas, def:spec}
    \lean{Specification.IsBrokenSymmetry}
    \leanok

    A symmetry $\tau$ of a specification $\gamma$ is \textbf{broken} if there is some
    $\mu \in \Gibbs{\gamma}$ with $\tau(\mu) \neq \mu$.
\end{definition}

\begin{theorem}[Georgii, the remark after (5.21): a broken symmetry is a phase transition]
    \label{thm:sym-broken-nontrivial}
    \uses{def:sym-broken, rmk:sym-transport-gibbs, cor:sym-unique-invariant}
    \lean{Specification.nontrivial_GP_of_isBrokenSymmetry,
          Specification.not_isBrokenSymmetry_of_GP_eq_singleton}
    \leanok

    If $\gamma$ has a broken symmetry then $|\Gibbs{\gamma}| > 1$. Contrapositively, if
    $\Gibbs{\gamma} = \set{\mu}$ then no symmetry of $\gamma$ is broken.
\end{theorem}
\begin{proof}
    \leanok

    Let $\tau$ be a symmetry with $\tau(\mu) \neq \mu$ for some $\mu \in \Gibbs{\gamma}$. By
    Remark \ref{rmk:sym-transport-gibbs}, $\tau(\mu) \in \Gibbs{\gamma}$ as well, so
    $\Gibbs{\gamma}$ contains two distinct elements. The contrapositive is
    Corollary \ref{cor:sym-unique-invariant}.
\end{proof}

\begin{example}[The two-dimensional Ising phase transition as a symmetry breaking]
    \label{ex:sym-ising-breaking}
    \uses{def:sym-broken, thm:sym-broken-nontrivial}
    \lean{MeasureTheory.GibbsMeasure.PeierlsSharp.isBrokenSymmetry_spinFlip,
          MeasureTheory.GibbsMeasure.PeierlsSharp.ising_symmetry_breaking,
          MeasureTheory.GibbsMeasure.PeierlsSharp.nontrivial_GP_of_symmetry_breaking}
    \leanok

    For the two-dimensional Ising ferromagnet in zero field with $\beta \ge \log 3$, the spin
    flip is a symmetry of the specification which is broken; consequently
    $|\mathcal{G}(\gamma)| > 1$.
\end{example}
\begin{proof}
    \leanok

    The Ising potential in zero field is spin-flip invariant, so the spin flip is a symmetry of
    the specification by Georgii (5.6)(c)/(5.9)(b). The Peierls estimate at
    $\log 9 \le 2\beta$ supplies two distinct Gibbs measures $\mu^+ \neq \mu^-$ exchanged by the
    spin flip, so $\mu^+$ is not spin-flip invariant and
    Theorem \ref{thm:sym-broken-nontrivial} applies.
\end{proof}
